\documentclass[a4paper,11pt]{article}
\usepackage[T2A]{fontenc}
\usepackage[utf8]{inputenc}
\usepackage[russian]{babel}
\usepackage{amsmath,amssymb,mathtools}
\usepackage{geometry}
\geometry{left=22mm,right=22mm,top=21mm,bottom=21mm}
\usepackage{microtype}
\usepackage{icomma}
\usepackage{tikz}
\usetikzlibrary{arrows.meta,positioning,shapes.geometric,calc}
\usepackage{booktabs,tabularx,array}
\usepackage{enumitem}
\usepackage{fancyhdr}
\usepackage{parskip}
\usepackage{xcolor}
\usepackage{hyperref}
\definecolor{accent}{HTML}{087D85}
\definecolor{darkaccent}{HTML}{075961}
\definecolor{textgray}{HTML}{34393C}
\definecolor{muted}{HTML}{646A6D}
\hypersetup{hidelinks,pdfauthor={Лёвин Артём Александрович},pdftitle={Чек-лист по алгебре. Володя Хачатурян. 17 ноября 2026 г.}}
\setlength{\parindent}{0pt}
\setlength{\parskip}{3.5pt}
\setlist[itemize]{leftmargin=16pt,itemsep=3pt,topsep=3pt}
\setlist[enumerate]{leftmargin=20pt,itemsep=5pt,topsep=3pt}
\renewcommand{\arraystretch}{1.25}
\pagestyle{fancy}
\fancyhf{}
\fancyhead[L]{\footnotesize\color{muted}Уравнения, неравенства и формулы}
\fancyhead[R]{\footnotesize\color{muted}Володя Хачатурян}
\fancyfoot[C]{\footnotesize\color{muted}\thepage}
\renewcommand{\headrulewidth}{0.35pt}
\renewcommand{\headrule}{\hbox to\headwidth{\color{accent}\leaders\hrule height \headrulewidth\hfill}}
\setlength{\headheight}{22pt}
\newcommand{\sectionline}{\par\vspace{1pt}\noindent\color{accent}\rule{\linewidth}{0.45pt}\color{textgray}\par\vspace{3pt}}
\newcommand{\minititle}[1]{\par\vspace{3pt}{\color{darkaccent}\bfseries #1}\par\vspace{1pt}}
\newcommand{\method}[1]{\textcolor{darkaccent}{\textbf{#1}}}
\newcommand{\topic}[1]{\section*{#1}\sectionline}
\newcommand{\sol}[1]{\[\displaystyle #1\]}
\newcommand{\ex}[2]{\begin{minipage}[t]{\linewidth}\textbf{#1.}\quad #2\end{minipage}\par\vspace{9pt}}
\setlength{\abovedisplayskip}{6pt}\setlength{\belowdisplayskip}{6pt}
\setlength{\abovedisplayshortskip}{4pt}\setlength{\belowdisplayshortskip}{4pt}
\AtBeginDocument{\color{textgray}}
\begin{document}

% ТИТУЛЬНЫЙ ЛИСТ
\thispagestyle{empty}
\begin{center}
\vspace*{22mm}
{\large\color{darkaccent}МАТЕМАТИКА\par}
\vspace{16mm}
{\LARGE Чек-лист\par}
\vspace{14mm}
{\Large\color{accent}\begin{tabular}{c}Уравнения, неравенства\\и формулы сокращённого умножения\end{tabular}\par}
\vspace{15mm}
{\normalsize Повторение алгебры: правила, примеры и самопроверка\par}
\vfill
{\large Володя Хачатурян\par}
\vspace{10mm}
{17 ноября 2026 г.\par}
\vspace{16mm}
{\small Лёвин Артём Александрович\\[3pt]
эксклюзивно для Володи Хачатуряна\par}
\vspace{15mm}
\begin{tikzpicture}[x=1cm,y=1cm]
\draw[accent,line width=0.85pt] (0,0) .. controls (1.1,.48) and (2.2,-.45) .. (3.5,0)
 .. controls (4.7,.48) and (5.8,-.45) .. (7,0);
\end{tikzpicture}
\vspace{6mm}
\end{center}
\clearpage
\setcounter{page}{1}

\topic{1. Деление десятичных дробей}
\method{Главное правило.} При делении десятичных дробей умножайте делимое и делитель на одну и ту же степень десяти так, чтобы в \emph{делителе} исчезла запятая. Если удобно, сделайте целыми оба числа.
\sol{1{,}4:0{,}25=\frac{1{,}4}{0{,}25}=\frac{140}{25}=5{,}6.}
\textbf{Почему можно так делать?} Значение частного не меняется, если числитель и знаменатель дроби умножить на одно и то же ненулевое число.
\minititle{Проверьте себя}
\begin{itemize}
  \item На сколько разрядов нужно перенести запятые?
  \item На одно и то же ли число умножены числитель и знаменатель?
  \item Можно ли проверить ответ обратным умножением?
\end{itemize}

\topic{2. Обратные действия в простых уравнениях}
\begin{center}
\begin{tabularx}{\linewidth}{@{}>{\raggedright\arraybackslash}X>{\raggedright\arraybackslash}X@{}}
\toprule
\textbf{Уравнение} & \textbf{Как найти неизвестное}\tabularnewline
\midrule
$a-b=c$ & $b=a-c$\\
$a:b=c$ & $b=a:c$, если $a\ne0$ и $c\ne0$\\
\bottomrule
\end{tabularx}
\end{center}
\textbf{Пример 1.} Решим $3-x=15$. Неизвестное вычитаемое равно разности уменьшаемого и результата:
\sol{x=3-15=-12.}
Проверка: $3-(-12)=15$.

\textbf{Пример 2.} Решим $\dfrac{3}{x}=15$. Здесь $x\ne0$; неизвестный делитель находим делением:
\sol{x=\frac{3}{15}=\frac15=0{,}2.}
Проверка: $3:0{,}2=15$.

\textbf{Пример 3.} Сначала раскрываем скобки, затем собираем подобные слагаемые:
\begin{align*}
-2(x+3)&=4(-x+1),\\
-2x-6&=-4x+4,\\
2x&=10,\\
x&=5.
\end{align*}
\method{Важно:} при переносе слагаемого в другую часть уравнения его знак меняется.

\topic{3. Квадратные уравнения и разложение на множители}
\minititle{Дискриминант}
Для уравнения $ax^2+bx+c=0$ ($a\ne0$):
\sol{D=b^2-4ac,\qquad x_{1,2}=\frac{-b\pm\sqrt D}{2a}.}
Если $D<0$, действительных корней нет. Если $D=0$, один корень. Если $D>0$, два различных корня.

\textbf{Пример.} $x^2-x-6=0$: здесь $a=1$, $b=-1$, $c=-6$.
\begin{align*}
D&=(-1)^2-4\cdot1\cdot(-6)=25,\\
x_1&=\frac{-(-1)+5}{2}=3,\qquad
x_2=\frac{-(-1)-5}{2}=-2.
\end{align*}
\method{Обязательно:} отрицательные коэффициенты при подстановке в формулу записывайте в скобках.

\topic{4. Сначала ищите общий множитель}
Перед вычислением дискриминанта задайте вопрос: \emph{можно ли вынести что-то за скобки?}
\sol{x^2-x=0\ \Longleftrightarrow\ x(x-1)=0.}
Произведение равно нулю, когда хотя бы один из множителей равен нулю:
\sol{x=0\quad\text{или}\quad x-1=0\ \Longrightarrow\ x=0,\;1.}

\minititle{Разность квадратов в уравнениях}
\begin{align*}
x^2-4=0
&\ \Longleftrightarrow\ (x-2)(x+2)=0
&&\Longrightarrow\ x=\pm2,\\
x^2-3=0
&\ \Longleftrightarrow\ (x-\sqrt3)(x+\sqrt3)=0
&&\Longrightarrow\ x=\pm\sqrt3.
\end{align*}
\textbf{Сравните:} $x^2+3=0$ не имеет действительных корней: $x^2\ge0$, поэтому $x^2+3>0$ при любом действительном $x$.

\minititle{Произведение нескольких множителей}
\sol{(x^2-10)(x^2-4x)=0.}
Приравниваем каждый множитель к нулю: $x^2-10=0$ или $x(x-4)=0$. Значит,
\sol{x\in\{-\sqrt{10},\;0,\;\sqrt{10},\;4\}.}
\method{Порядок действий:} сначала вынесение общего множителя и формулы сокращённого умножения; дискриминант --- для квадратичных множителей, которые не удалось решить проще.

\topic{5. Линейные неравенства}
\minititle{Строгие и нестрогие знаки}
\begin{center}
\begin{tabularx}{\linewidth}{@{}p{31mm}XX@{}}
\toprule
 & \textbf{На числовой прямой} & \textbf{Граничная скобка}\\
\midrule
$x>a$ или $x<a$ & выколотая точка & круглая: $(\,)$\\
$x\ge a$ или $x\le a$ & закрашенная точка & квадратная: $[\,]$\\
\bottomrule
\end{tabularx}
\end{center}
\textbf{Пример.} $x-1{,}2\ge0$ означает $x\ge1{,}2$, то есть $x\in[1{,}2;+\infty)$.

\textbf{Пример с переносом и сменой знака:}
\begin{align*}
2(x-3)&\le3(x-5),\\
2x-6&\le3x-15,\\
-x&\le-9,\\
x&\ge9.
\end{align*}
На последнем шаге обе части умножены на $-1$, поэтому знак неравенства \textbf{обращён}.
\begin{center}
\begin{tikzpicture}[x=1cm,y=1cm]
\draw[-{Stealth[length=2mm]},gray!70] (0,0)--(9,0) node[right]{$x$};
\draw[accent,very thick,-{Stealth[length=2.5mm]}] (4,0)--(8.3,0);
\fill[accent] (4,0) circle (2.2pt);
\draw[muted] (4,-.13)--(4,.13);
\node[below=5pt] at (4,0) {$9$};
\node[above=6pt,accent] at (6.5,0) {$x\ge9$};
\end{tikzpicture}
\end{center}
\textbf{Ответ:} $x\in[9;+\infty)$.
\textbf{Запомните:} при умножении или делении \emph{обеих частей неравенства} на отрицательное число знак меняется на противоположный.

\topic{6. Системы неравенств: ищем пересечение}
Система означает, что \emph{все} её условия должны выполняться одновременно.
\sol{\begin{cases}x\ge1{,}5,\\x<10\end{cases}
\quad\Longrightarrow\quad x\in[1{,}5;10).}
\begin{center}
\begin{tikzpicture}[x=1cm,y=1cm]
\foreach \h in {0,0.65} {\draw[-{Stealth[length=2mm]},gray!65] (0,\h)--(9,\h);}
\draw[accent,thick,-{Stealth[length=2.4mm]}] (2,.65)--(8.45,.65);
\fill[accent] (2,.65) circle (2.1pt);
\draw[accent,thick,-{Stealth[length=2.4mm]}] (7,0)--(.5,0);
\fill[white] (7,0) circle (2.5pt);\draw[accent,line width=.9pt] (7,0) circle (2.5pt);
\draw[darkaccent,line width=2.2pt] (2,-.42)--(7,-.42);
\fill[darkaccent] (2,-.42) circle (2.1pt);
\fill[white] (7,-.42) circle (2.5pt);\draw[darkaccent] (7,-.42) circle (2.5pt);
\node[above=6pt] at (2,.65) {$1{,}5$};
\node[above=6pt] at (7,.65) {$10$};
\node[left] at (0,.65) {\small $x\ge1{,}5$};
\node[left] at (0,0) {\small $x<10$};
\node[below=9pt,darkaccent] at (4.5,-.42) {\small Общее решение};
\end{tikzpicture}
\end{center}
Границы записываются в порядке возрастания; $1{,}5$ входит в ответ, а $10$ не входит.

\textbf{Пример посложнее:}
\begin{align*}
\begin{cases}
2(x-3)+1\ge -x+5,\\
3-3(x+2)\le4x-7
\end{cases}
&\Longleftrightarrow
\begin{cases}3x\ge10,\\7x\ge4\end{cases}\\[4pt]
&\Longleftrightarrow
\begin{cases}x\ge\dfrac{10}{3},\\x\ge\dfrac47.\end{cases}
\end{align*}
Обе штриховки направлены вправо, поэтому общая часть начинается с большего числа $\frac{10}{3}$:
\sol{\boxed{x\in\left[\frac{10}{3};+\infty\right).}}
\textbf{Важное правило:} решение каждого неравенства по отдельности --- промежуточный результат. Для системы всегда нужно найти \emph{пересечение} решений.

\clearpage
\topic{7. Формулы сокращённого умножения}
\minititle{Разность квадратов}
\sol{\boxed{a^2-b^2=(a-b)(a+b).}}
Это формула для \emph{разности двух квадратов}. Она позволяет разложить выражение на множители, а также свернуть произведение обратно в разность.
\begin{align*}
4a^2-b^2&=(2a-b)(2a+b),\\
16x^2-9y^2&=(4x-3y)(4x+3y),\\
(3x-y)(3x+y)&=9x^2-y^2.
\end{align*}
\textbf{Ограничение:} $a^2+b^2$ нельзя разложить по формуле разности квадратов: между квадратами должен быть \emph{минус}.

\minititle{Квадрат разности и квадрат суммы}
\sol{\boxed{(a-b)^2=a^2-2ab+b^2,}}
\sol{\boxed{(a+b)^2=a^2+2ab+b^2.}}
В этих формулах $a$ и $b$ --- целые выражения; при подстановке удобно брать их в скобки.

\textbf{Пример 1.} В $(2x-y)^2$ вместо $a$ подставляем $2x$, вместо $b$ --- $y$:
\begin{align*}
(2x-y)^2
&=(2x)^2-2\cdot(2x)\cdot y+y^2\\
&=4x^2-4xy+y^2.
\end{align*}
\textbf{Пример 2.} При раскрытии $(5t-3k)^2$:
\begin{align*}
(5t-3k)^2
&=(5t)^2-2\cdot(5t)\cdot(3k)+(3k)^2\\
&=25t^2-30tk+9k^2.
\end{align*}
\textbf{Пример 3.} Если $(4k-4t)^2$, среднее слагаемое равно $-2\cdot4k\cdot4t=-32kt$:
\sol{(4k-4t)^2=16k^2-32kt+16t^2.}
\method{Критично:} минус перед $b$ уже содержится в формуле $-2ab$. Величину $b$ подставляйте \emph{без дополнительного минуса}; последний член $+b^2$ всегда имеет знак «плюс».

\minititle{Сравните выражения}
Разность квадратов и квадрат разности дают разные результаты:
\begin{align*}
(4k)^2-(3t)^2&=(4k-3t)(4k+3t),\\
(4k-3t)^2&=16k^2-24kt+9t^2.
\end{align*}

% АЛГОРИТМ НА ОТДЕЛЬНОЙ СТРАНИЦЕ
\clearpage
\begin{center}
{\Large\color{darkaccent}Алгоритм: система линейных неравенств}\par
\vspace{3mm}
\textcolor{muted}{Последовательность действий при решении системы}\par
\vspace{10mm}
\begin{tikzpicture}[
  node distance=9mm,
  every node/.style={font=\small,align=center},
  startstop/.style={ellipse,draw=darkaccent,line width=.75pt,minimum width=46mm,minimum height=11mm},
  stage/.style={rectangle,rounded corners=2pt,draw=darkaccent,line width=.75pt,text width=91mm,minimum height=15mm,inner sep=5pt},
  decision/.style={diamond,aspect=2.4,draw=darkaccent,line width=.75pt,text width=37mm,inner sep=2.8pt},
  finish/.style={rectangle,rounded corners=2pt,draw=accent,line width=.85pt,text width=48mm,minimum height=17mm,inner sep=5pt},
  conn/.style={-{Stealth[length=2.5mm]},draw=accent,line width=.9pt}]
\node[startstop] (s) {Дана система неравенств};
\node[stage,below=of s] (a) {Решите \textbf{каждое} неравенство отдельно.\\При делении на отрицательное число измените знак.};
\node[stage,below=of a] (b) {Изобразите оба решения на числовых прямых.\\Проверьте граничные точки.};
\node[decision,below=11mm of b] (d) {Есть общая часть\,?};
\node[finish,below=17mm of d] (yes) {Запишите пересечение решений в виде промежутка};
\node[finish,right=20mm of d] (no) {Общих значений нет.\\Ответ: $\varnothing$};
\draw[conn] (s)--(a);
\draw[conn] (a)--(b);
\draw[conn] (b)--(d);
\draw[conn] (d)--node[right,near start,fill=white,inner sep=1pt]{\footnotesize да}(yes);
\draw[conn] (d.east)--node[above,near start,fill=white,inner sep=1pt]{\footnotesize нет}(no.west);
\end{tikzpicture}
\par\vspace{10mm}
{\small\textbf{Итог:} ответ системы должен удовлетворять одновременно всем неравенствам.\par}
\end{center}

% ТРЕНИРОВКА
\clearpage
\section*{Тренировочный блок}
\sectionline
\textbf{Путь А --- 10 заданий.} Работайте последовательно. В заданиях на неравенства укажите ответ промежутком; при необходимости изображайте числовую прямую. Ответы можно сверить с таблицей на следующей странице.
\vspace{3mm}
\minititle{Базовый уровень}
\ex{1}{Вычислите: $1{,}8:0{,}3$.}
\ex{2}{Решите уравнение: $7-x=19$.}
\ex{3}{Решите уравнение: $\dfrac6x=15$.}
\minititle{Средний уровень}
\ex{4}{Решите квадратное уравнение: $x^2-x-12=0$.}
\ex{5}{Решите уравнение: $x^2-16=0$.}
\ex{6}{Решите неравенство: $-2x+4\ge10$.}
\ex{7}{Разложите на множители: $25t^2-16k^2$. Укажите название формулы.}
\ex{8}{Раскройте скобки: $(3a-2b)^2$. Укажите название формулы.}
\minititle{Уровень выше среднего}
\ex{9}{Решите систему неравенств и запишите ответ промежутком:
\[\begin{cases}2(x-2)+1\ge -x+3,\\5-(x+2)>0.\end{cases}\]}
\ex{10}{Решите уравнение, используя разложение на множители:
\[(x^2-4)(x^2-5x+6)=0.\]}

% ОТВЕТЫ
\clearpage
\section*{Ответы к тренировочному блоку}
\sectionline
\vspace{4mm}
\begin{center}
\begin{tabularx}{.86\linewidth}{@{}>{\centering\arraybackslash}p{13mm}X@{}}
\toprule
\textbf{№} & \textbf{Ответ}\tabularnewline
\midrule
1 & $6$\\
2 & $x=-12$\\
3 & $x=\dfrac25$, $x\ne0$\\
4 & $x=-3$ или $x=4$\\
5 & $x=-4$ или $x=4$\\
6 & $x\in(-\infty;-3]$\\
7 & $(5t-4k)(5t+4k)$ --- разность квадратов\\
8 & $9a^2-12ab+4b^2$ --- квадрат разности\\
9 & $x\in[2;3)$\\
10 & $x\in\{-2;2;3\}$\\
\bottomrule
\end{tabularx}
\end{center}
\vspace{14mm}
\minititle{Как проверить себя}
\begin{itemize}
\item Подставьте найденные корни в исходные уравнения.
\item Для неравенств проверьте одно число внутри и одно число вне найденного промежутка.
\item Разложения и раскрытия скобок проверьте обратным умножением.
\item При записи системы проверьте, что все её условия выполняются одновременно.
\end{itemize}
\vfill
\begin{center}
\small\color{muted}Лёвин Артём Александрович \textbullet\ эксклюзивно для Володи Хачатуряна\\
17 ноября 2026 г.
\end{center}
\end{document}
